\section{The closure and the hull of the feasible set}\label{app:feasible-hull}

This appendix proves the statements of Section~\ref{sec:feasible-hull} about
the closure $\mathcal C$ and the set $\Sigma$.

\begin{proposition}[Free affine columns]\label{prop:free-remainders}
Write $v=(x,z)$ with $x$ the block, and let $M$ be the matrix formed by the
columns of $A$ that belong to $z$, with one row for each inequality and one
row for each equality (of the two opposite rows of an equality, $M$ keeps one).
If $M$ has full row rank, then $\mathcal C=\conv\Sigma$.
\end{proposition}

\begin{proof}
Let $v=(x,z)\in\mathcal C$. By Carath\'eodory's theorem there are $u_j\in D$ and
weights $\theta_j>0$ summing to one with $\sum_j\theta_ju_j=x$ and
$\bar y:=\sum_j\theta_jg(u_j)\le b-Av$, with equality in the rows of each
equality. Let $A_x$ denote the block columns of $A$, restricted like $M$, let
$M^+$ be a right inverse of $M$, put $\rho_j=\bar y-g(u_j)+A_x(x-u_j)$ in the
same rows, and $z_j=z+M^+\rho_j$. A direct substitution shows that each row
of $(u_j,z_j)$ has the same left-hand side as $\bar y+Av$, so
$(u_j,z_j)\in\Sigma$. Since $\sum_j\theta_j\rho_j=0$, the convex combination of
the points $(u_j,z_j)$ is $v$.
\end{proof}

The objective row always contains $-t$, so for the objective alone the
closure is the hull of the epigraph over $D$. The setting of
\citet{LiersEtAl2021}, rows $z_j=g_j(x)$ with a separate variable $z_j$ for
each row, is also covered by Proposition~\ref{prop:free-remainders}. Without
free affine columns the gap can be strict.

\begin{example}[Gap to the feasible hull]\label{ex:quarter}
Let $D=[0,1]$ and take the equality row $x^2=1/4$, whose feasible set is
$\{1/2\}$. By Corollary~\ref{cor:one-row}, $\mathcal C=\{x:x^2\le1/4\le x\}=[1/4,1/2]$,
and the two cuts $x\ge1/4$ and $x\le1/2$ describe it. The point $x=1/4$ is the
mean of the measure with mass $3/4$ at $0$ and $1/4$ at $1$, which satisfies
the row in expectation.
\end{example}

This is a Lagrangian duality gap \citep{Geoffrion1974,Nowak2005}: the rows are
imposed on averages, after convexification. Even the complete family of cuts
accepts $x=1/4$, so the gap is not a failure of direction search. It is
closed by convexifying aggregated sets instead of aggregated functions, as in
the aggregation closure of \citet{DeyMunozSerrano2022}: here
$\conv\{x\in[0,1]:x^2\le1/4\}=[0,1/2]$ and $\conv\{x\in[0,1]:x^2\ge1/4\}=[1/2,1]$
intersect in $\{1/2\}$. It is also closed by placing the row inside $D$, which
changes the support problem; our implementation keeps nonlinear rows out of
$D$ so that the support problems stay in the classes of
Section~\ref{sec:quadratic}. With
$\Sigma_\lambda=\{v:\ Sv\in D,\ \lambda\T(g(Sv)+Av-b)\le0\}$ for $\lambda\ge0$, the
hierarchy is
\[
\conv\Sigma\ \subseteq\ \bigcap_{\lambda\ge0}\conv\Sigma_\lambda\ \subseteq\ \mathcal C .
\]
The middle set is the convexified closure of the surrogate relaxations of
\citet{MullerEtAl2022}, and the gap to $\mathcal C$ is a surrogate--Lagrangian
gap. Example~\ref{ex:quarter} shows that the second inclusion can be strict.
The first can be strict as well, even when $D$ is exactly the projection of
$\Sigma$: with $D=[0,2]$, a free $z$ and the rows $z\le(x-2)^2/2$ and
$z\le x^2/2$, the set $\Sigma$ satisfies $z\le\min\{x/2,1-x/2\}$ on its hull,
so $z\le1/2$ at $x=1$. Every aggregated row reads $z\le q_\lambda(x)$ with
$q_\lambda=[\lambda_1(x-2)^2+\lambda_2x^2]/(2(\lambda_1+\lambda_2))$ convex and
$q_\lambda(0)+q_\lambda(2)=2$, so the chord gives $(1,1)\in\conv\Sigma_\lambda$ for
every $\lambda$, and $(1,1)$, the mean of the point masses at $x=0$ and $x=2$,
also lies in $\mathcal C$.
